% ADD/REMOVE THE 'answers' OPTION TO INCLUDE/SUPPRESS SOLUTIONS
\documentclass[11pt,addpoints]{exam}
% \documentclass[11pt,addpoints,answers]{exam}

\newcommand{\hwnum}{4}
\newcommand{\duedate}{September 30}

% In order to compile this file you will need to get 'header.tex'
% and make the line below point to the appropriate file path
\RequirePackage{microtype}
\RequirePackage{mathtools}
\RequirePackage{amsthm}
\RequirePackage{amssymb}
\RequirePackage{xspace}
\RequirePackage[shortlabels]{enumitem}
\RequirePackage{xcolor}
\RequirePackage{hyperref}
\RequirePackage[capitalize,nameinlink,noabbrev]{cleveref} % must load after hyperref
\usepackage{fvextra} % To use Verbatim
\usepackage{diagbox} % For diagbox
\RequirePackage[boxed]{algorithm}
\RequirePackage[noend]{algpseudocode}
\RequirePackage{tikz}
\usetikzlibrary{arrows,automata,positioning}

\hypersetup{breaklinks=true,
    colorlinks=true,
    linkcolor=blue,
    filecolor=blue,
    citecolor=blue,
    urlcolor=blue}

\algrenewcommand{\algorithmiccomment}[1]{\texttt{//} #1}
\algrenewcommand\algorithmicrequire{\textbf{Input:}}
\algrenewcommand\algorithmicensure{\textbf{Output:}}


% allow cleveref to label and reference enumerables defined in the exam class.
% these automatically define corresponding \Crefname as well.

% https://tex.stackexchange.com/questions/126020/cleveref-doesnt-use-correct-capitalized-name-if-used-with-amsthm
\makeatletter
\if@cref@capitalise
\crefname{question}{Question}{Questions}
\Crefname{partno}{Part}{Parts}
\crefname{subpart}{Subpart}{Subparts)}
\crefname{subsubpart}{Subsubpart}{Subsubparts}
\else
\crefname{question}{question}{questions}
\crefname{partno}{part}{parts}
\crefname{subpart}{subpart}{subparts}
\crefname{subsubpart}{subsubpart}{subsubparts}
\fi
\makeatother

% numeric sets in "blackboard" font

\newcommand{\N}{\mathbb{N}}
\newcommand{\Z}{\mathbb{Z}}
\newcommand{\R}{\mathbb{R}}
\newcommand{\Q}{\mathbb{Q}}
\newcommand{\C}{\mathbb{C}}
\newcommand{\Prop}{\mathbb{P}}

% paired delimiters

\DeclarePairedDelimiter\abs{\lvert}{\rvert}
\DeclarePairedDelimiter\length{\lVert}{\rVert}
\DeclarePairedDelimiter\norm{\lVert}{\rVert}
\DeclarePairedDelimiter\parens{(}{)}
\DeclarePairedDelimiter\tuple{(}{)}
\DeclarePairedDelimiter\brackets{[}{]}
\DeclarePairedDelimiter\floor{\lfloor}{\rfloor}
\DeclarePairedDelimiter\ceil{\lceil}{\rceil}
\DeclarePairedDelimiter\round{\lfloor}{\rceil}
\DeclarePairedDelimiter\set{\{}{\}}
\DeclarePairedDelimiter\inner{\langle}{\rangle}

\newcommand{\bit}{\set{0,1}}

% asymptotics

\DeclareMathOperator{\Otil}{\tilde{O}}
\DeclareMathOperator{\poly}{poly}
\DeclareMathOperator{\polylog}{polylog}
\DeclareMathOperator{\negl}{negl}

% algorithms

\newcommand{\algo}[1]{\textsc{#1}}

\newcommand{\memo}{\text{memo}}
\newcommand{\tabl}{\text{table}}
\newcommand{\backtrack}{\text{backtrack}}

\newcommand{\ALG}{\text{ALG}}
\newcommand{\OPT}{\text{OPT}}
\newcommand{\weight}{\text{weight}}
\newcommand{\val}{\text{value}}

% computability

% named language
\newcommand{\lang}[1]{L_{\text{#1}}}
% computational problem
\newcommand{\cproblem}[1]{\ensuremath{\text{#1}}\xspace}
% class of languages
\newcommand{\class}[1]{\ensuremath{\mathsf{#1}}\xspace}

\newcommand{\qst}{q_{\text{start}}}
\newcommand{\qacc}{q_{\text{acc}}}
\newcommand{\qrej}{q_{\text{rej}}}

\newcommand{\Lbarber}{\lang{BARBER}}
\newcommand{\atm}{\lang{ACC}}
\newcommand{\htm}{\lang{HALT}}
\newcommand{\ehtm}{\lang{$\varepsilon$-HALT}}
\newcommand{\eqtm}{\lang{EQ}}
\newcommand{\etm}{\lang{$\emptyset$}}
\newcommand{\epstm}{\lang{$\set{\varepsilon}$}}
\newcommand{\Lprop}{\lang{$\Prop$}}
\newcommand{\LSigmastar}{\lang{$\Sigma^*$}}

% complexity

\newcommand{\yes}{\ensuremath{\text{YES}}}
\newcommand{\no}{\ensuremath{\text{NO}}}

\newcommand{\DTIME}{\class{DTIME}}
\renewcommand{\P}{\class{P}}
\newcommand{\NP}{\class{NP}}
\newcommand{\NPH}{\class{NPH}}
\newcommand{\NPC}{\class{NPC}}
\newcommand{\coNP}{\class{coNP}}

\newcommand{\MAZE}{\cproblem{MAZE}}
\newcommand{\PALINDROME}{\cproblem{PALINDROME}}
\newcommand{\TSP}{\cproblem{TSP}}
\newcommand{\SAT}{\cproblem{SAT}}
\newcommand{\CSAT}{\cproblem{CSAT}}
\newcommand{\TSAT}{\cproblem{3SAT}}
\newcommand{\VC}{\cproblem{VERTEX-COVER}}
\newcommand{\DS}{\cproblem{DOMINATING-SET}}
\newcommand{\SC}{\cproblem{SET-COVER}}
\newcommand{\HC}{\cproblem{HAMCYCLE}}
\newcommand{\HP}{\cproblem{HAMPATH}}
\newcommand{\AHP}{\cproblem{ANYHAMPATH}}
\newcommand{\IS}{\cproblem{IS}}
\newcommand{\CLIQUE}{\cproblem{CLIQUE}}
\newcommand{\SSUM}{\cproblem{SUBSET-SUM}}
\newcommand{\KNAPSACK}{\cproblem{KNAPSACK}}
\newcommand{\MAXCUT}{\cproblem{MAX-CUT}}

% randomness

\DeclareMathOperator*{\Var}{Var}
\DeclareMathOperator*{\Ex}{\mathbb{E}}

\newcommand{\RP}{\class{RP}}
\newcommand{\coRP}{\class{coRP}}
\newcommand{\BPP}{\class{BPP}}
\newcommand{\ZPP}{\class{ZPP}}
\newcommand{\BQP}{\class{BQP}}

%%% misc

\newcommand{\eps}{\varepsilon}

%%% theorems

\theoremstyle{plain}            % following are "theorem" style

\newtheorem{theorem}{Theorem}
\newtheorem{lemma}[theorem]{Lemma}
\newtheorem{corollary}[theorem]{Corollary}
\newtheorem{proposition}[theorem]{Proposition}
\newtheorem{claim}[theorem]{Claim}
\newtheorem{fact}[theorem]{Fact}
\newtheorem{openproblem}[theorem]{Open Problem}

\theoremstyle{definition}       % following are def style

\newtheorem{definition}[theorem]{Definition}
\newtheorem{conjecture}[theorem]{Conjecture}
\newtheorem{protocol}[theorem]{Protocol}
\newtheorem{exercise}[theorem]{Exercise}

\theoremstyle{remark}           % following are remark style

\newtheorem{example}[theorem]{Example}
\newtheorem{remark}[theorem]{Remark}
\newtheorem{note}[theorem]{Note}

%%% for homework and section notes

\newcommand{\commonheader}[2]{
    \pagestyle{headandfoot}
    \setlength{\headheight}{26pt}
    \setlength{\headsep}{16pt}

    \header
        {\small{\textbf{EECS 376: Foundations of Computer Science}} \\ \footnotesize{\textbf{University of Michigan, Fall 2026}}}
        {#1}
        {#2}

    \firstpageheadrule
    \runningheadrule

    \footer
        {}
        {\thepage}
        {}
}

\newcommand{\hwheader}{
    \commonheader
        {\Large \textbf{Homework \hwnum}}
        {\small \textbf{Due 8:00pm, \duedate\\ {\tiny(accepted until 10:00 pm, no credit after)}}}
}

\newcommand{\hwslnheader}{
    \commonheader
    	{}
        {\Large \textbf{Solutions to Homework \hwnum}}
    \printanswers
}

\newcommand{\notesheader}{
    \commonheader
    	{}
        {\Large \textbf{Discussion Notes \sectionnum}}
}

\newcommand{\practiceheader}{
    \commonheader
    	{}
        {\Large \textbf{Discussion Worksheet \sectionnum}}
}

\newcommand{\practiceslnheader}{
    \commonheader
    	{}
        {\Large \textbf{Solutions to Discussion Worksheet \sectionnum}}
}

\newcommand{\reviewheader}{
    \commonheader 
    \smallskip
    	{}
        {\Large \textbf{Midterm Review Notes}}
}

\newcommand{\hwpreface}{

\noindent This homework has \numquestions\ questions, for a total of \numpoints\ points and \numbonuspoints\ extra-credit points.

\noindent Unless otherwise stated, each question requires \emph{clear}, \emph{logically correct}, and \emph{sufficient} justification to convince the reader.

\noindent We strongly encourage you to typeset your solutions in \LaTeX.

\noindent If you collaborated with someone, you must state their name(s). You must \emph{write your own solution} for all problems and \emph{may not use any other student’s write-up or any AI tools}.
}

\newcommand{\hint}[1]{
\emph{Hint}: #1
}

% exam class setup
\pointsinmargin
\pointpoints{pt}{pts}
\bonuspointpoints{EC pt}{EC pts}
\marginpointname{ \points}
\marginbonuspointname{ \bonuspoints}

\usepackage{tikz}
\usetikzlibrary{arrows.meta,positioning}

\ifprintanswers
\hwslnheader   % header for solutions
\else
\hwheader   % header for homework
\fi

\begin{document}


\hwpreface

\begin{questions}

  \addtocounter{question}{-1}
  \question \textbf{Before you start; before you submit.}

      

    If applicable, state the name(s) and uniqname(s) of your collaborator(s).

    \begin{solution}
       
    \end{solution}


\question[10] \textbf{Self assessment.} \nopagebreak
  Carefully read and understand the posted solutions to the previous homework.
  Identify one part for which your own solution has the most room for improvement (e.g., has unsound reasoning, doesn’t show what was required, could be significantly clearer or better organized, etc.).
  Copy or screenshot this solution, then in a few sentences, explain what was deficient and how it could be fixed.

  (Alternatively, if you think one of your solutions is significantly \emph{better} than the posted one, copy it here and explain why you think it is better.)

  If you didn't do last week's homework, choose a problem from it that looks challenging to you, and in a few sentences, explain the key ideas behind its solution in your own words.

  \begin{solution}

  
  
        % Easter egg for those using LaTeX!
        % Here's how you can insert an image in LaTeX, in case you need to include a screenshot
        % \includegraphics[width=0.5\linewidth]{waquackquack.jpg} % Adjust the dimension as you need
        
  \end{solution}

\question \textbf{Fixed-color-sequence paths.} 
% \jrnote{Graph DP question from W24}
% \fxnote{In the future say the graph representation is adjacency list}
    
 
  Let $G=(V,E)$ be an undirected graph with~$n$ vertices and~$m$ edges, where $m \geq n$.
  Also, suppose that each edge is colored either black or white.
  We say that a path $P=(e_{1},\dots,e_{k})$ of edges has \emph{color sequence} $(c_{1},\dots,c_{k})$, where each $c_{i}\in\set{\text{black},\text{white}}$, if each $e_{i}$ has color $c_{i}$.
      In this problem, you will design an algorithm that, given two vertices~$s,t$ and a desired color sequence $(c_{1},\dots,c_{k})$, determines in $O(mk)$ time whether there exists a path of length $k$ from~$s$ to~$t$ with that color sequence.
  Note that paths in this problem do not need to be \emph{simple}; they can visit vertices or edges more than once.

  The graph $G=(V,E)$ is represented via \emph{adjacency lists}, i.e., for each vertex $i \in V$ there is a list of all its neighbors: vertices~$j$ such that $(i,j) \in E$.

  \begin{parts}
    \part[10]
    Derive a recurrence relation, including base case(s), for this problem.
    Briefly justify its correctness.
      
    \begin{solution}
    \end{solution}
        
    \part[6] Give pseudocode implementing an $O(km)$-time dynamic programming solution to this problem, and justify its running time.
    You do not need to give a correctness analysis since you did so for your recurrence. 
    
    \begin{solution}
    \end{solution}
  \end{parts}

  \question \textbf{Back to college} 
  \nopagebreak
   Michael drives down to Ann Arbor from his hometown for the start of each school year.
   Interested in minimizing his commute time, he invents a modular fuel pack for his car for the commute through various towns in Michigan.
   This fuel pack allows the vehicle to travel a certain distance before refueling.
   Each town has a service station where a fuel pack can be reloaded: the current pack is simply swapped out for a fresh full one.
   (Any remaining fuel in the removed pack will not be used by the car.)
   The goal is to visit the towns in a specified order, while minimizing the number of fuel pack reloads.

   We model this problem as follows.
   The car can travel a distance up to~$R$ on a single fuel pack.
   Starting from town 0, there are~$n$ other towns to visit in sequence.
   For $i=1, \ldots, n$, the distance from the $(i-1)$st town to the $i$th town is $d_i \leq R$.
   Given~$R$ and an array $D = [d_1, \ldots, d_n]$ as input, the goal is to find a \emph{smallest} set $S \subseteq \set{1,\ldots,n}$ of towns (specified by their indices) where refueling should occur, so that the car can reach all~$n$ towns without running out of fuel at any point.
   The car starts out with a full fuel pack, and no refueling is needed at the final town.

  \begin{parts}
    \part[7] Give a greedy algorithm (including pseuodocode) that solves this problem.
    Defer a correctness argument to the next part, but provide everything else involved in ``giving an algorithm'' here.
    
    \begin{solution}
    \end{solution}
    
    \part[10] The setup for your correctness proof is as follows.
    Let $\OPT$ be some arbitrary optimal solution.
    Let $\ALG=i_1,i_2,i_3,\dots, i_k$ be the towns chosen by your algorithm, in their greedily chosen order.
    If $\ALG = \OPT$, then $\ALG$ is optimal, as needed.
    So suppose that $\ALG \neq \OPT$, and consider the \emph{first} town $i_{\text{diff}}$ in $\ALG$ that is not in $\OPT$.
    (Such an $i_{\text{diff}}$ must exist because $\ALG$ must have at least as many towns as $\OPT$, since $\OPT$ is optimal.)
    So, $\OPT$ contains all of $i_1,i_2,i_3,\dots, i_{\text{diff}-1}$, but not $i_{\text{diff}}$.

    You will construct another optimal solution $\OPT'$ that contains all of $i_1,i_2,i_3,\ldots,i_{\text{diff}}$ by modifying $\OPT$, exchanging out some suitable town~$j$ for $i_{\text{diff}}$ instead.

    Describe which town~$j$ to exchange for $i_\text{diff}$, and prove that $\OPT'$ is an optimal solution. In particular, you should prove that $\OPT'$ is a valid solution \emph{and} an optimal solution. 
    
    \begin{solution}
    \end{solution}

    \part[4] Explain in your own words why the previous part (proving that $\OPT'$ is an optimal solution that contains $i_1, \ldots, i_\text{diff}$) ultimately proves that your algorithm is correct.

    \begin{solution}
    \end{solution}
    
    \part[5] Now, suppose that the price of a fuel pack may differ from town to town.
    For example, a fuel pack might cost \$8 in the first town, \$3 in the second town, etc. So, the input now also includes an array $C=[c_1, \ldots, c_n]$, where~$c_i$ is the cost of refueling in the $i$th town.
    The goal now is to complete the tour while \emph{minimizing the total cost} of the fuel packs that are purchased along the way.

    Michael claims that we should use the same greedy algorithm that you proposed above to solve this problem.
    Determine whether this claim is true or not.
    \begin{itemize}
        \item If it is true, give a proof using a greedy exchange argument.
        \item Otherwise, provide a small counterexample: give a specific input, the output of your greedy algorithm on that input, and a better solution for that input; also, \emph{very briefly} explain where your proof from the previous part fails for this problem. 
    \end{itemize}
    

    \begin{solution}
    \end{solution}
  \end{parts}

  
  \question \textbf{Dynamic programming for dynamic shortest paths.} \nopagebreak
             
  Graphs that arise in real applications are often not fixed, but can change over time; such graphs are called ``dynamic.''
  (This is a totally different use of the word ``dynamic'' than in ``dynamic programming''.
  Computer scientists sometimes name things in confusing ways!)
  For example, roads and intersections can be added to or deleted from the road network, computers can be added to or removed from the Internet, and people can (un)follow each other on social networks.
  For many problems of interest, there are algorithms for dynamic graphs that are faster than just re-computing answers ``from scratch'' whenever the graph changes.
  In this problem you will give such algorithms for shortest-path problems.

  \vspace{1.0cm}
  
  Let $G=(V, E)$ be a weighted directed graph with $n=\abs{V}$ vertices, where the weight of each edge $(u,v)$ is denoted $w(u,v)$.
  (There can be negative-weight edges, but there is no negative-weight cycle in $G$.)
  Suppose that we have already computed the all-pairs distance table~$D$ of~$G$.
  That is, for each vertex pair $u,v \in V$, entry $D(u,v)$ stores the distance (i.e., the length of a shortest path) from~$u$ to~$v$ in $G$.

  Now, a new vertex $v_{new}$ is added to the graph, together with its incident edges.
  Let $G_{new} = (V \cup \set{v_{new}}, E \cup E_{new})$ denote the updated graph, where $E_{new}$ consists of all the incoming and outgoing edges for $v_{new}$.
  (There is no negative-weight cycle in $G_{new}$.)
  We aim to compute the updated distance table $D_{new}$ of~$G_{new}$.

  A straightforward approach is to compute $D_{new}$ from scratch using the Floyd-Warshall algorithm, in $\Theta(n^3)$ time.
  But we want to do better by exploiting the fact that we already have~$D$.
  In this problem, you will obtain a faster $O(n^2)$-time algorithm.

  \begin{parts}
    
    \part[10] Write expressions for $D_{new}(v_{new},u)$ and $D_{new}(u,v_{new})$ that hold for all $u\in V$, where the expressions are in terms of the already-known quantities $D(y,z)$ for $y,z\in V$, and $w(y,z)$ for $(y,z) \in E \cup E_{new}$. 
    Evaluating your expressions should take $O(n)$ time for each vertex~$u$, for a total of $O(n^2)$ time.
    Justify the correctness of your expressions and the evaluation time.
    
    \hint{Consider why the Bellman-Ford algorithm is correct.}
    
    \begin{solution}
    \end{solution}

    \part[8] Write an expression for $D_{new}(u,v)$ that holds for all $u,v \in V$, where the expression is in terms of the already-known quantities listed in the previous part, as well as the quantities you computed in the previous part.
    Evaluating your expression for all $u,v \in V$ should take a total of $O(n^2)$ time.
    Justify the correctness of your expression and the evaluation time.

    Observe that by combining the two parts, we have computed the entire new distance table $D_{new}$ of $G_{new}$ in $O(n^2)$ time!
    
    \hint{Consider why the Floyd-Warshall algorithm is correct.}

    \begin{solution}
    \end{solution}
  \end{parts}  



  \question \textbf{MST mashup.} \nopagebreak 
  
  \begin{parts}
    
    \part[5] Consider the following claim: a minimum spanning tree in a graph is \emph{unique} if the edge weights are \emph{distinct} (all different).

    Professor $\Upsilon$ presents you with the following bogus ``proof'' of this claim: \emph{because edge weights are distinct, all choices in Kruskal's algorithm are completely determined---there are no ``ties'' between edge weights that the algorithm can choose how to break.
      So, the algorithm has only one possible output.
      Because Kruskal's algorithm is correct for the MST problem, its unique output must be the only minimum spanning tree in the graph.}
      
    Briefly explain what is wrong with this ``proof.''
    In particular, indicate which sentence in the ``proof'' does not follow from the previous one, and provide a reason for it.

    \begin{solution}
    \end{solution}

    \part[8]
    Give a correct proof for the claim in the previous part.
    
    \hint{Consider two distinct spanning trees $T_1, T_2$ and prove via an exchange argument that one of them does not have minimum weight.}

    \begin{solution}
    \end{solution}

  \end{parts}


  \question \textbf{Greedy knapsack filling.} \nopagebreak
    
  Recall the knapsack problem (sometimes called the ``0-1'' knapsack problem): you are given~$n$ items, where the $i$th item has weight~$w_{i}$ and value~$v_{i}$ (both positive integers), and a non-negative weight capacity~$C$ of the knapsack.
  An optimal solution is a subset $S\subseteq \set{1, 2, \ldots, n}$ of the items having maximum total value \[ \sum_{i\in S} v_i \; \text, \] under the constraint that the total weight is at most the knapsack capacity, i.e., \[ \sum_{i\in S} w_i \leq C \; \text. 
  \]

  In this problem, we introduce the \emph{fractional} variant of the knapsack problem, where one may take any fraction $p_i \in [0,1]$ of any item~$i$.
  Naturally, a $p_i$-fraction of item~$i$ weighs $p_i \cdot w_i$, and has value $p_i \cdot v_i$.
  The goal is to maximize the total value of the selected fractional items.

  \begin{parts}

    \part[2] 
    Briefly explain why the optimal value for the fractional knapsack problem is \emph{at least as large} as that of the original 0-1 knapsack problem (for the same weights, values, and knapsack capacity).

    \begin{solution}
    \end{solution}
    
    \part[5] Consider the following two greedy algorithms for the fractional knapsack problem:
    \begin{itemize}
    \item \algo{ValueGreedy}: While there is still some unused knapsack capacity and at least one unconsidered item, choose an item having the \emph{largest value} (among those not considered yet), and add the largest possible fraction of that item (up to the entire item) that will fit within the remaining knapsack capacity.
        
    \item \algo{WeightGreedy}: While there is still some unused knapsack capacity and at least one unconsidered item, choose an item having the \emph{smallest weight} (among those not considered yet), and add the largest possible fraction of that item (up to the entire item) that will fit within the remaining knapsack capacity.
    \end{itemize}

    Suppose we run \emph{both} algorithms \algo{ValueGreedy} and \algo{WeightGreedy}, and output a best one of their two outputs.
    Is this a correct algorithm for the fractional knapsack problem?
    If so, prove it.
    Otherwise, give an input for which this algorithm's output is not optimal, and show why it is not.

    \begin{solution}
    \end{solution}

    \part[10] Consider the following greedy algorithm for the fractional knapsack problem: 
    \begin{itemize}
    \item \algo{RelativelyGreedy}: While there is still some unused knapsack capacity, choose an item having the largest \emph{relative value} $v_i/w_i$ (among those not considered yet), and add the largest possible fraction of that item (up to the full item) that will fit within the remaining knapsack capacity.
    \end{itemize}

    In this part, you will prove that the \algo{RelativelyGreedy} algorithm is a correct algorithm for the fractional knapsack problem.
    This is notable because we used a more complicated dynamic programming algorithm to solve 0-1 Knapsack, but the fractional variant admits a much simpler greedy algorithm!

    The setup for your correctness proof is as follows.
    Without loss of generality, suppose that the items are in sorted order by relative value $r_{i}=v_{i}/w_{i}$, from largest to smallest (the algorithm considers them in this order).
    Let $\ALG=p_1,p_2,\dots,p_n$ be the fractions of items $1,2,\dots, n$ chosen by the algorithm, respectively.
    (So, $p_i=1$ if the algorithm chooses all of item~$i$, and $p_i=0$ if it doesn't choose any of item~$i$.)

    Let $\OPT$ be an arbitrary optimal solution.
    If $\ALG = \OPT$, then $\ALG$ is optimal, as needed.
    So suppose that $\ALG \neq \OPT$, and consider the smallest index $j$ such that $\OPT$ does not have \emph{exactly} a $p_j$-fraction of item $j$.

    You will construct another solution $\OPT'$ whose fractions of the first~$j$ items are exactly $p_1,p_2,\ldots,p_j$ by modifying $\OPT$, using an exchange argument.

    Describe which fraction(s) of item(s) to exchange, and prove that $\OPT'$ is an optimal solution; this should include a proof that $\OPT'$ is a \emph{valid} solution.

    \begin{solution}
    \end{solution}

  \end{parts}
  

  \bonusquestion \textbf{Optional extra credit: Greedily independent.} 

  
  \nopagebreak

  An \emph{independent set} of an undirected graph $G=(V,E)$ is a subset of vertices $S \subseteq V$ such that no two vertices in $S$ are adjacent, i.e., $(u,v) \notin E$ for every distinct $u,v \in S$.
  We are interested in a \emph{maximum} independent set, i.e., the independent set with the \emph{largest} number of vertices.
  (There may be more than one maximum independent set.)
    
  Consider the following greedy algorithm for finding an independent set in a graph.

  \begin{minipage}{\linewidth}
    \begin{algorithm}[H]
      \begin{algorithmic}[1]
        \Function{GreedyIS}{$G=(V,E)$}
        \State{$S \gets \emptyset$}
        \While{$G$ has at least one vertex}
            \State{Select any vertex~$v$ in $G$ that has \emph{smallest} degree (i.e., the fewest neighbors)}
            \State{Add $v$ to $S$}
            \State{Remove $v$ and all its neighbors, including all their incident edges, from $G$}
        \EndWhile
        \State{\Return{$S$}}
        \EndFunction
      \end{algorithmic}
    \end{algorithm}
  \end{minipage}

  \begin{parts}

    \bonuspart[1]
    
    Give an example of a graph for which the algorithm \algo{GreedyIS} \textbf{does not} return the maximum independent set.
    You can assign numbers to the vertices of your graph, and assume that in the case where there are multiple vertices with the smallest degree, \algo{GreedyIS} will choose the vertex having the smallest number. 
    \begin{solution}
    \end{solution}

    
    \bonuspart[2] In this part, you will prove that \algo{GreedyIS} returns a maximum independent set when the input graph is a \emph{tree}.

    The setup for the proof is as follows.
    Let $\ALG=v_1,v_2,v_3,\dots,v_k$ be the vertices chosen by your algorithm, in their greedily chosen order, and let $\OPT$ be some arbitrary optimal solution.
    If $\ALG = \OPT$, then $\ALG$ is optimal, as needed.
    So suppose that $\ALG \neq \OPT$, and consider the \emph{first} vertex $v_{\text{diff}}$ in $\ALG$ that is not in $\OPT$.
    (Such a $v_{\text{diff}}$ must exist because $\ALG$ cannot be a proper subset of $\OPT$; do you see why?)
    So, $\OPT$ contains all of $v_1,v_2,v_3,\dots,v_{\text{diff}-1}$, but not $v_{\text{diff}}$.
    
    You will construct another optimal solution (a maximum independent set) $\OPT'$ that has all of $v_1,v_2,v_3,\ldots,v_{\text{diff}}$ by modifying $\OPT$, exchanging out some suitable vertex~$u$ for $v_{\text{diff}}$ instead.

    Describe which vertex~$u$ to exchange for $v_\text{diff}$, and prove that $\OPT'$ is both valid \textit{and} optimal.

    \begin{solution}
    \end{solution}

  \end{parts}




\end{questions}
\end{document}